\subsection*{Circle-to-Grid-Maze Collision Detection}

\textbf{Prompt:}
\begin{lstlisting}[language={}, breaklines=true]
I am building a 2D tilt-controlled game called "BiteBound" for an ESP32-S3 paired with an ST7789 display (240x240px). The gameplay consists of a ball roll navigating a custom maze grid. We represent our board as a 1D uint8_t flat array of size width * height, where 0 represents a free corridor and 1 represents a wall.

I need a high-performance C++ class called `MazeCollider` that checks collisions for a circular `PhysicsBody` (holding x, y, and radius). To keep frame rates high on the microcontroller, the check should only scan wall cells within the body's bounding box instead of checking the whole board. It should return a `Contact` struct containing:
- `hit` (bool)
- `penetration` (float, maximum intersection depth)
- `normal` (Vec2, indicating the direction to push the body out of the wall)

Could you write the execution method for this collider? Keep the physics-to-grid mapping simple (1x1 units per cell) and optimize for the ESP32's floating-point unit.
\end{lstlisting}

\textbf{Answer:}
\begin{lstlisting}[language={}, breaklines=true]
// Algorithm: Bounding-Box Grid-Circle Collision Detection
function queryCollision(body):
    r = body.radius
    
    // 1. BOUNDING BOX SPATIAL PRUNING
    // Limit grid checks strictly to tiles overlapping the circle's bounding box
    x0 = floor(body.x - r)
    x1 = ceil(body.x + r)
    y0 = floor(body.y - r)
    y1 = ceil(body.y + r)
    
    // Clamp coordinates to remain within grid boundaries
    x0 = max(0, x0)
    y0 = max(0, y0)
    x1 = min(grid_width - 1, x1)
    y1 = min(grid_height - 1, y1)

    maxPenetration = 0.0
    collisionNormal = [0.0, -1.0] // Default fallback normal
    hit = false

    // 2. TILE EVALUATION
    for py from y0 to y1:
        for px from x0 to x1:
            if board[px, py] is not WALL:
                continue // Skip empty corridors
                
            // 3. CLOSEST POINT PROJECTON (AABB Clamping)
            // Clamp circle center onto the 1x1 tile bounds
            nx = clamp(body.x, px, px + 1)
            ny = clamp(body.y, py, py + 1)
            
            // 4. METRIC COMPUTATIONS
            dx = body.x - nx
            dy = body.y - ny
            distanceSquared = dx*dx + dy*dy
            
            // 5. NARROW-PHASE SATISFIABILITY TEST
            if distanceSquared < r * r:
                distance = sqrt(distanceSquared)
                penetration = r - distance
                
                // Track only the deepest cell penetration for single resolver
                if penetration > maxPenetration:
                    maxPenetration = penetration
                    hit = true
                    
                    if distance > 1e-4:
                        collisionNormal = [dx / distance, dy / distance]
                    else:
                        collisionNormal = [0.0, -1.0] // Perfect overlap fallback

    if hit:
        return Contact(hit=true, penetration=maxPenetration, normal=collisionNormal)
    else:
        return Contact(hit=false, penetration=0.0, normal=[0.0, 0.0])
\end{lstlisting}

\subsubsection*{Collision Logic & Structural Breakdown}

The algorithm implements a **static Circle-vs-AABB (Axis-Aligned Bounding Box)** collision pipeline specialized for a uniform 2D grid:

\begin{enumerate}
    \item \textbf{AABB Spatial Hashing/Narrow-phase Pruning:} Instead of checking the ball against all pixels in the entire maze, it determines the outer limits of the ball in grid coordinates ($x_0, x_1, y_0, y_1$) using local flooring and ceiling operations. It limits evaluation to a maximum of a $3\times3$ or $4\times4$ grid area around the body, bringing the complexity down from $O(W \times H)$ to $O(R^2)$ (where $R$ is the body's radius relative to cell size).
    \item \textbf{Grid-cell Closest Point Projection:} For each colliding grid cell, it calculates the closest point $(nx, ny)$ inside that specific $1\times1$ pixel block to the circle's center $(body.x, body.y)$. This is achieved by clamping the circle center coordinates to the cell's physical bounding boundaries.
    \item \textbf{Narrow-Phase Separation Test:} It computes the vector from the closest cell point to the circle center. If its magnitude is less than the radius, there is a physical overlap.
    \item \textbf{Local Deepest Penetration Filter:} Rather than accumulating all overlapping impulses, the loop selectively records only the cell that yields the deepest penetration depth ($maxPen$). The contact normal is configured to point directly from that closest point to the center, indicating the shortest path vector required to push the ball out of the wall.
\end{enumerate}

\subsubsection*{Typical Misconceptions \& Vulnerabilities}

When developers utilize or adapt this code, they often encounter several physical and mathematical misconceptions:

\begin{itemize}
    \item \textbf{The Single-Deepest-Penetration Fallacy (Corner Stuckness):} 
    By only tracking the single deepest penetration, the algorithm ignores multi-cell overlaps. If the ball rolls into an internal corner (colliding with both a bottom wall and a right wall), only the deeper wall is resolved. In the next frame, the physics engine pushes the ball directly sideways into the secondary wall block, leading to visual stuttering, deceleration, or the player getting stuck. A full impulse solver should accumulate response forces from all intersecting cells rather than selecting just one.
    
    \item \textbf{Perfect Overlap Singularity (Zero-Distance Division):} 
    If the ball's center lands exactly on a cell border or perfectly inside a wall cell, the distance $d \approx 0$. At this point, normal calculation ($dx/d, dy/d$) becomes undefined or undergoes division by zero. The fallback normal `Vec2(0.0f, -1.0f)` is arbitrarily hardcoded. If the actual exit path was upwards, this fallback pushes the ball deeper into the wall, causing it to clip through.
    
    \item \textbf{The Tunneling Vulnerability (High-Speed Clip):} 
    This algorithm performs static intersection tests. If the ball is moving at high speeds (where its displacement per frame exceeds its radius), it can jump over wall pixels completely between frames. To resolve this, continuous collision detection (CCD) or sub-stepping within the physics loop is needed.
    
    \item \textbf{Uniform Scale Assumption:} 
    This engine assumes a hardcoded scale where 1 cell = 1 unit of physical coordinate. If the rendering or physical world scales up (e.g., each tile is $16\times16$ display pixels), this method will break unless coordinates are transformed into local grid space before the bounding box calculation.
\end{itemize}